\chapter{Вимірювання вологості у виробничому приміщенні}

\section{Виробниче приміщення для дослідження}

Для проведення дослідів було обрано виробниче приміщення з трьома вікнами, дверним отвором та одним
робочим місцем. План приміщення зображений на \ref{fig:building}. Контрольні зони в яких проводились
заміри зазначені цифрами 1---5.

\begin{figure}[h] \centering
  \includegraphics[width=0.9\textwidth]{./images/building.svg.png}
  \caption{План виробничого приміщення.}
  \label{fig:building}
\end{figure}

\section{Результати спостереження}

\begin{table}[ht]
\caption{Результати спостереження}
\label{t:results}
\begin{tabular}{| c | c | c | c | c | c |}
\hline
\multicolumn{1}{|p{2cm}|}{КЗ,~№} &
\multicolumn{1}{C{2.2cm}|}{1} &
\multicolumn{1}{C{2.2cm}|}{2} &
\multicolumn{1}{C{2.2cm}|}{3} &
\multicolumn{1}{C{2.2cm}|}{4} &
\multicolumn{1}{C{2.2cm}|}{5} \\ \hline 
Вологість,~\% & 1,51 & 1,44 & 1,85 & 1,11 & 1,22 \\ \hline 
\end{tabular}
\end{table} 


